\section{Gestione dei file}
Il file di dump di un database è tipicamente un file di grosse dimensioni, il quale ha al suo interno diverse tabelle che riportano le credenziali degli utenti, che sono state criptate con qualche algoritmo di hashing. Dall'altra parte sarà necessario un file, più o meno grande, contenente le credenziali più usate che si vuole utilizzare per l'attacco. Lo scambio e la gestione di questi due file all'interno del cluster ha un suo costo, sopratutto perchè è spesso un trade off tra la porzione di RAM che possiamo occupare e il carico sull'I/O dovuta alla lettura intensiva dei file. Supponendo infatti di ragionare per estremi, caricando completamente entrambi i file in memoria oppure di leggerlo dal filesystem ogni qualvolta ve ne sia bisogno si può subito evincere che, nel caso di file di dimensioni importanti, nel primo caso è molto probabile che si occupi tutto lo spazio in RAM disponibile, nel secondo le unità di esecuzione saturerebbero così velocemente la coda delle richieste dell'I/O che avremmo un throughput pessimo del processore, con conseguente spreco di tempo sul cluster, che tipicamente comporta un costo economico irragionevole. 

\subsection{Equa suddivisione}
Partendo da un metodo piuttosto semplice, si consideri il caso i cui la rete sia composta da un manager e P worker, il manager decide a che porzioni di dati i worker devono accedere facendo si che siano in tutte equali. La suddivisione può avvenire in due modi, all'inizio il manager assegna ad ogni nodo una porzione di database da attaccare, oppure dinamicamente un nodo chiede al manager una porzione di database da testare quando finisce quella precedente. \\
Considerando il caso statico, il file viene diviso e trasmesso prima che i workers possano cominciare il loro task, per cui considerando K bytes da condividere bisognerà trasmettere $(K/P)*P$ bytes e ogni worker dovrà allocare $K/P$ bytes di memoria. In qualche modo significa comunque saturare molto l'I/O nelle prime fasi di esecuzione, senza parallelamente svolgere alcuna computazione, sacrificando quindi tutti i vantaggi di avere a bordo dell'hardware specializzato nel gestire asincronamente la memoria, come ad esempio i DMA e i NIC dei singoli nodi, che svolgono le loro attività indipendentemente dallo stato del processore. Va da se che lo scenario si ripropone anche in caso di un filesystem clusterizzato, infatti tutti i worker all'inizio proveranno ad accedere al file contenuto nell'unità remota, la quale farà attendere i worker mentre trasmette i dati, sprecando tempo di esecuzione. 
\\ 
Il caso dinamico è invece anche a livello concettuale, visibilmente migliore. Le sezioni del file vengono decise a priori dal manager, ma i worker nè richiedono la trasmissione solo quando finiscono la precedente. All'inizio della computazione 
